% ==========================================
% 8. DISKUSSION [cite: 86, 215]
% ==========================================
\section*{Diskussion}
\addcontentsline{toc}{section}{Diskussion}
%Här presenterar ni er tolkning av resultaten och bedömer deras signifikans. Diskutera möjliga konsekvenser av resultaten, och presentera eventuella rekommendationer. Det är viktigt att diskussionen använder resultaten för att besvara er frågeställning.

%Avsnittet ska också innehålla reflektioner kring arbetet, som till exempel dess begränsningar. Ni kan också diskutera lösningar på problem som ni identifierat och diskuterat tidigare, eller ta upp andra problem som arbetet inte behandlat, frågor som ej besvarats.

%Koppla eventuellt också dina resultat till tidigare arbeten. På så sätt kan diskussionen bli ett samtal med det ni skrev i tidigare avsnitt. Slutligen ska ni sätta ditt eget arbete i ett större sammanhang, bredda ditt perspektiv. Kan dina resultatet generaliseras? Kan det du gjort användas i något annat sammanhang?


Vi kom fram till slut på ett resultat från experimenten som ger tydliga svar på rapportens frågeställning om huruvida latenta representationer från SDXL uppvisar bildlik egenskaper som kan optimeras med traditionell komprimering utan betydande kvalitetsförlust. Genom att analysera effektivitet och begränsningar


\subsection*{Tolkning av resultat }
Vi upptäckte att JPEG XL (lossless) presterar bättre än det generella NPZ. Att JPEG XL uppnår ett lägre BPP-värde (0,413 mot 0,417) vid identisk bildkvalitet indikerar att datan i det latenta rummet inte är stokastisk. Detta bekräftar att SDXL:s latenta representationer bibehåller en strukturell ordning som är tillräckligt hög för att kunna behandlas med traditionella bildkodningsmetoder.Dock bör det noteras att skillnaden mellan JXL lossless och NPZ är marginell, vilket tyder på att även generella komprimeringsalgoritmer för binärdata är relativt effektiva på detta område.
En upptäckt i studien är att kvantisering fungerar väldigt väl. Genom att reducera precisionen från 32-bitars flyttal till 8-bitars heltal (RGBA) bevisar vi att den latenta representationen är robust. Decodern lyckas återskapa bilden med hög kvalitet trots denna förenkling, vilket är en av arbetets starkaste slutsatser.
Vid förlust komprimering lossy uppvisar resultaten en betydande robusthet vid låga distansvärden (d=0,05 och d=0,1). Trots en reduktion i PSNR på upp till 1,26 dB bibehålls en hög visuell trohet. Detta tyder på att VAE decodern kan tolerera en viss grad av kvantiseringsbrus i latenten utan att den slutgiltiga rekonstruktionen försämras. Denna observation är signifikant, då den möjliggör en reduktion av filstorleken med cirka 16 % utan att kompromissa med bildens semantiska integritet.

\subsection*{Begränsningar och konsekvenser}
En tydlig begränsning som identifierats är den kraftiga kvalitetsförluster vid aggressivare komprimering. Vid d=0,5 och d=1,0 sjunker PSNR och SSIM drastiskt. Detta tyder på att de latenta särdragen är mer känsliga för förstörande komprimering än vad vanliga RGB pixlar ofta är medan latenten tål att vi rundar av värden genom kvantisering, tål den inte att information gallras bort genom aggressivare lossy-algoritmer.
En annan aspekt att beakta är spridningen i resultat mellan olika motiv. Bild 13 uppvisade exempelvis betydligt lägre PSNR (19,82 dB) än bild 3 (31,60 dB). Detta innebär att metoden inte är universellt förutsägbar utan beror kraftigt på bildens innehåll.
Arbetet visar att det finns en stor potential i att lagra bilddata i ett ML-format (maskininlärningsvänligt format). Istället för att lagra tunga RGB-filer som måste avkodas innan de matas in i en modell, kan man lagra kvantiserade latenta representationer direkt. Detta kan generaliseras till andra områden, såsom videoströmning för autonoma system eller storskalig lagring av medicinska bilder där både mänsklig granskning och maskinell analys krävs. Genom att använda SDXL:s format med 4 kanaler (RGBA) möjliggörs dessutom användning av befintlig infrastruktur för bildhantering, vilket sänker tröskeln för implementering.
Framtida forskning bör utforska om nyare modeller som FLUX uppvisar samma känslighet för lossy-komprimering, eller om andra bitdjup än 8 bitar skulle kunna utgöra en mer optimal balansgång för lagring av latenta särdrag.

